% ============================================================
% 03_AI-reflection.tex
% Section 3: Building with AI
% ============================================================

\section{Building with AI}\label{sec:building}

\subsection{Human Judgment in Managing AI Output}

The primary challenge in developing Maverick-SORT was not generating code, but verifying its accuracy, feasibility, and relevance at the sheer speed of AI output. While the agent (Claude Code) translated my mathematical KGDRL models into PyTorch boilerplate efficiently, it suffered from severe context degradation: it frequently forgot macro-project flows, struggled with cross-branch GitHub push logistics, and failed to retrieve previously working commands efficiently. I also found it challenging when trying to maintain one practical solution when the agent itself was searching aimlessly so I had to interrupt the running and redirect it back to the right way.  

To counter this, the human role shifted to strict \textbf{project management and traceability}. AI-generated outputs were only accepted when they satisfied three criteria: \textbf{\textcolor{red}{technical correctness, consistency with the KGDRL design, and practical relevance to the warehouse use case}}. I forced the agent to maintain continuous markdown logs---such as \texttt{ai\_usage\_review.md}, \texttt{research\_note.md}, and \\ \texttt{DATA\_UPLOAD\_FEASIBILITY\_REPORT.md}---recording every dependency change, architectural pivot, and run result. When the AI inevitably hallucinated or drifted off-course, these logs allowed me to rapidly audit its steps, find the causes, and explicitly guide it back to the established R\&D baseline. AI is a revolutionary means of production, but human agency and project orchestration remain the absolute driving forces of creation. The resulting division of labor was simple: the AI accelerated execution, while I remained responsible for direction, validation, and integration.

% \subsection{A Concrete Error: Systems Thinking vs. Local Patching}

% One instructive failure during the Streamlit Cloud deployment perfectly illustrates the limitation of the AI's local reasoning compared to human systems-level thinking. 

% The AI had generated a config file with a standard \texttt{import transformers}. On the cloud, this triggered a catastrophic cascade: the import scanned for image models, sought \texttt{torchvision} (which was absent), threw a \texttt{ModuleNotFoundError}, and inadvertently triggered Streamlit's background file-watcher to infinitely restart the app.

% When asked to fix this, the AI proposed adding \texttt{torchvision==0.22.0} to the requirements. However, this conflicted with our pinned \texttt{torch==2.12.0}. The AI then blindly suggested bumping \texttt{torch}, which would have destroyed our CUDA compatibility. 

% \textbf{Human judgment}: The AI was treating the symptom, not the disease. Recognizing the root cause, I implemented a systems-level defense: (1) suppressing the scan via environment variables, (2) lazy-loading the import, and crucially, (3) disabling Streamlit's file watcher entirely\parencite{Jurowetzki2026}. The AI lacks the environmental awareness to understand cloud background-watcher dynamics; that operational wisdom must come from the human.

\subsection{A Concrete Error: Version Iteration and Context Amnesia}

One instructive failure occurred not in localized code logic, but in version iteration management, perfectly illustrating the AI's severe context amnesia across a longitudinal project. 

When generating successive versions of the application (e.g., iterating through multiple \texttt{streamlit\_app.py} versions), the AI consistently failed to maintain feature parity. Instead of executing strict incremental updates based on the project blueprint, the AI produced fractured ``updates'' plagued by hallucinations and arbitrarily dropped functionalities. The outputs were often regressions rather than true iterations. 

\textbf{Human judgment}: I initially expected the AI to seamlessly carry forward previous capabilities, failing to establish a stable, automated iteration mechanism to constrain it. Consequently, development became highly inefficient. After every AI generation, I had to expend massive cognitive bandwidth manually auditing the new code against the \texttt{ai\_usage\_review.md} logs and the original plan to detect silent deletions. The AI treats each iteration as an isolated event; human judgment and rigorous manual QA were strictly required to enforce longitudinal consistency.

\subsection{Defensive Engineering and the ``Happy Path'' Fallacy}

Similarly, we found that \textbf{AI intrinsically assumes a ``happy path.''} For the CSV upload module, the AI wrote ingestion code expecting the user's data to perfectly match our schema. Knowing the hostility and unpredictability of actual warehouse data, I manually intervened to enforce defensive programming (e.g., \texttt{df.columns.intersection(expected)}). The AI can generate the processing logic, but only the human understands the chaos of the operational environment it will be deployed in, shifting the paradigm from \textit{``assume the data is right''} to \textit{``assume the data is broken.''}

% --- AI Reflection 逻辑总结图 ---
\begin{figure}[htbp]
\centering
\begin{tikzpicture}[
    node distance=0.6cm and 0.5cm,
    box/.style={draw=primarycolor, thick, rounded corners, align=center, inner sep=6pt, fill=primarycolor!5, text width=4.5cm, minimum height=2.8cm},
    conclusion/.style={draw=accentcolor, thick, rounded corners, align=center, inner sep=8pt, fill=accentcolor!10, text width=14.5cm},
    arrow/.style={->, thick, draw=aaugray, >=stealth}
]

% Pillar 1: 3.1
\node[box] (b1) {
    \textbf{3.1 Human Judgment}\\[0.2cm]
    \footnotesize \textbf{Challenge:}\\ Validation $>$ Generation\\[0.1cm]
    \footnotesize \textbf{Process:}\\ Markdown Traceability\\[0.1cm]
    \footnotesize \textbf{Filter:}\\ ``Good Enough'' Criteria
};

% % Pillar 2: 3.2
% \node[box, right=of b1] (b2) {
%     \textbf{3.2 Systems Thinking}\\[0.2cm]
%     \footnotesize \textbf{Context:}\\ Cloud Deployment Crash\\[0.1cm]
%     \footnotesize \textbf{AI Flaw:}\\ Local Patching (Symptom)\\[0.1cm]
%     \footnotesize \textbf{Human Fix:}\\ Root Cause Defense
% };

% Pillar 2: 3.2
\node[box, right=of b1] (b2) {
    \textbf{3.2 Version Iteration}\\[0.2cm]
    \footnotesize \textbf{Context:}\\ App Version Updates\\[0.1cm]
    \footnotesize \textbf{AI Flaw:}\\ Context Amnesia \& Drops\\[0.1cm]
    \footnotesize \textbf{Human Fix:}\\ Manual Audit \& QA
};

% Pillar 3: 3.3
\node[box, right=of b2] (b3) {
    \textbf{3.3 Defensive Eng.}\\[0.2cm]
    \footnotesize \textbf{Context:}\\ CSV Upload Module\\[0.1cm]
    \footnotesize \textbf{AI Flaw:}\\ The ``Happy Path'' Fallacy\\[0.1cm]
    \footnotesize \textbf{Human Fix:}\\ Defensive Integration
};

% Conclusion Block
\node[conclusion, below=1.2cm of b2] (conc) {
    \textbf{\large The Ultimate Division of Labor}\\[0.15cm]
    \textit{AI accelerates execution; humans provide direction, validation, and integration.}
};

% Connecting Arrows
\draw[arrow] (b1.south) -- (b1.south |- conc.north);
\draw[arrow] (b2.south) -- (conc.north);
\draw[arrow] (b3.south) -- (b3.south |- conc.north);

\end{tikzpicture}
\caption{Logical framework of Human-AI collaboration during system development.}
\label{fig:ai_reflection_summary}
\end{figure}
% --- 图表结束 ---